
\section{算法与Infra——不同优化层的一体两面}

\subsection{核心洞察：从问题出发，反向推导}

视频开场即点出了一个被许多人纠结的问题：算法工程师还是Infra工程师？几位嘉宾的核心观点出奇一致——\textbf{这本质上是同一个问题的不同优化层面}。

\begin{quote}
从上层来讲，模型或算法的优化，一直延伸到下层的系统优化，这些东西其实都是相通的。关键还是看你要解决什么样的问题，然后从这个问题的角度去``walking backwards''，去想不同的优化算法和优化角度。
\end{quote}

\begin{importantbox}{核心主张：优化是连续的，不是对立的}
算法优化和系统优化位于同一个优化频谱的两端，而非两个互斥的领域。真正的区分不在于你做哪一端，而在于你从哪个入口切入问题。无论从算法层往下走，还是从系统层往上走，最终目标都是端到端地解决一个真实问题。
\end{importantbox}

\subsection{MLC：算法与系统交汇的典范}

讨论中引用了机器学习编译（MLC，Machine Learning Compilation）作为典型案例。MLC 要求从业者同时具备两个维度的能力：

\begin{itemize}
    \item \textbf{ML 知识}：理解模型的数学结构、计算图、算子语义
    \item \textbf{System 知识}：掌握编译器后端、内存管理、硬件适配
\end{itemize}

没有 ML 背景的纯系统工程师很难理解为什么某个算子需要特殊的编译优化；没有系统背景的纯算法工程师则无法将优化落地到实际硬件上。MLC 的成功恰恰说明，\textbf{在两个领域的交界处，往往蕴藏着最大的创新机会}。

\subsection{看看顶尖人物的轨迹}

嘉宾提到，观察行业内顶尖人物的职业轨迹是一个很有价值的方向判断方法：

\begin{knowledgebox}{参考坐标}
李牧（Mu Li）、蒋家青（Jiaqing Jiang）、庞若明（Ruoming Pang）等——这些在 ML 和 System 都有标志性贡献的研究者，从职业生涯早期就展现出算法与 Infra 的双重能力，既有算法 Paper，也有 Infra Paper。这并非偶然，而是由问题本质决定的。
\end{knowledgebox}

\subsection{本章小结}

算法和 Infra 不是两个互斥的职业选择，而是同一问题在不同抽象层级上的优化。MLC 是这一理念的产物，而行业顶尖人物的双重能力则是最好的旁证。关键不是选 A 还是 B，而是从你要解决的问题出发，沿着优化链条向下或向上延伸自己的能力边界。

\vspace{1em}
\begin{figure}[H]
\centering
\includegraphics[width=0.78\textwidth,keepaspectratio]{figures/talk_algo_infra.jpg}
\caption{几位嘉宾在圆桌讨论中分享观点 \protect\footnotemark}
\end{figure}
\footnotetext{视频画面时间区间：00:08:20--00:10:40。}


\section{职业选择：产品还是Infra？}

\subsection{两种路径的本质差异}

讨论将职业路径分为两个大方向：靠近产品的团队和靠近 Infra 的团队。两者各有清晰的利弊：

\begin{itemize}
    \item \textbf{产品侧}：曝光度高，直接接触用户，能获得第一手信息和市场反馈。如果产品成功，个人价值可能获得十倍、百倍的增长。但产品失败后，积累的经验不一定可迁移。
    \item \textbf{Infra 侧}：工作更``苦''，曝光度低，不太站在聚光灯下。但技术积累具有\textbf{线性甚至超线性的增长特性}——日复一日的深耕不会因为某个产品的成败而清零。
\end{itemize}

\begin{importantbox}{Infra 的长期复利}
Infra 能力的积累类似于复利曲线：它可能起步慢、初期回报低，但随着时间推移，你在系统设计、性能优化、故障排查等方面的经验会形成不可替代的护城河。产品方向的回报则更像期权——大多数产品失败，但少数成功带来的回报是指数级的。
\end{importantbox}

\subsection{招聘市场的残酷现实}

嘉宾分享了两组极具冲击力的数字：

\begin{itemize}
    \item \textbf{算法岗位}：一个岗位可能收到 \textbf{1:10000} 的简历比例
    \item \textbf{Infra 岗位}：通常只收到 \textbf{20--50 份}简历
\end{itemize}

这个差距并非短期现象，而是``持续存在，甚至还在扩大化''。算法方向的热度吸引了大量涌入者，而 Infra 方向因为门槛高、见效慢，人才供给始终相对稀缺。

\subsection{创业视角：端到端能力是底线}

从创业者的角度，嘉宾给出了一个直白到近乎残酷的判断：

\begin{quote}
作为一个创过业的人，我百分之百不会招那种只有算法能力的算法同学。在真的创业的时候，更重要的是服务客户，而不是追求什么 fancy 的算法。大多数算法都能从网上的开源拿到，你只要通过数据和各种方式把它变成产品。如果你只是把 Transformer 从 HuggingFace 拉下来做成一个模型，但没有能力把它部署到 AWS 上，那我觉得这种同学的价值非常有限。
\end{quote}

\begin{warningbox}{不要做``丢包侠''}
有一种风气：算法同学做出模型后直接丢给 Infra 团队说``你们去 scale 吧''，自己不 end-to-end 地 own 整个流程。职业初期也许可以这样，但到了中后期，这种工作方式会严重限制你的 skill set 和 credit 积累。你做的算法，你必须能把它部署出去。
\end{warningbox}

\subsection{本章小结}

产品侧和 Infra 侧各有围城，没有绝对的好坏。但市场数据清楚地表明：Infra 方向竞争更小、积累更稳；产品方向则回报更高但风险也更大。无论是哪条路，端到端的交付能力——从算法设计到生产部署——都是不可或缺的底线。

\vspace{1em}
\noindent 本节内容来自视频字幕时间区间：00:03:00--00:10:00。


\section{PhD 方向选择：信噪比与时机判断}

\subsection{最危险的策略：追当前最火的方向}

PhD 是五年时间窗口的长周期决策。嘉宾对如何选择研究方向给出了一个重要判断原则：

\begin{importantbox}{赌方向的第一原则}
\textbf{你不能做现在最火的东西，你要做未来可能会火的东西。}如果你做现在最火的东西，PhD 毕业时这个方向可能已经凉了，或者已经饱和了——所有人都在里面竞争。
\end{importantbox}

具体案例：Object Detection。几年前是 CV 领域最火的方向，招一个 Scientist 能收到几百份简历。现在呢？``人均 Faster R-CNN Detector Two''，MASK R-CNN 遍地都是，发一篇 Object Detection 的文章可能连 WACV 都不一定能中。

\subsection{``隔壁王阿姨''信号：已变现，不是信号}

嘉宾用了一个极具画面感的比喻来说明信号与噪声的区别：

\begin{quote}
当你家隔壁的王阿姨都开始跟你讲 DeepSeek 的时候，其实就证明你已经不该搞了。已经凉了。该上船的人都上船了，该占的坑位都已经被占走了。
\end{quote}

这不仅仅是一个玩笑。它的深层含义是：当信息穿透到完全没有专业壁垒的圈层时，所有的先发优势和时间窗口都已经关闭。这时候你听到的不是信号，而是\textbf{噪声}——或者更准确地说，是已经变现完毕的过期信号。

\begin{warningbox}{信号 vs 已变现信息}
不要用行业顶尖 10--20 个人的行为来判断方向。那个级别的信息已经不是信号了，已经是完全的过去式，已经被他们变现了。你应该看的是：哪些方向正在被少数聪明人悄悄布局，但还没有变成大众话题。
\end{warningbox}

\subsection{PhD 中期 Pivot 是可行的}

嘉宾也给了希望：PhD 方向不是锁死的。很多人是在博三、博四时突然把方向转到 LLM 方向，在剩下的一两年内往死里干，出来几篇 Paper，也走出了不错的出路。关键不是一开始就选对，而是保持\textbf{足够的灵活性和执行力}。

\subsection{地理信息差：硅谷比新加坡早一个节拍}

一个被反复提及的点：地理 proximity 对方向判断的影响。

\begin{quote}
在硅谷，GPT-3 的时候已经很多人开始嗅到这个味道了。但在新加坡，一直到 GPT-3.5 之后才会嗅到这个味道。
\end{quote}

所以嘉宾建议：你有条件的话，每年去硅谷、深圳这些前沿地区待一段时间，找人聊一聊。这种面对面获取的``软信息''，比论文和新闻早 6--12 个月。

\subsection{不超过5年的判断}

最后一层原则：

\begin{importantbox}{思考的时间尺度}
不要做超过 5 年以上的思考和决定。这个时代变化太快，技术在以跃迁式的方式发展。你只能在周期之内做判断——比如互联网的 2000--2008，移动互联网的 2008--2018。接下来的 3 年，你就赌你相信的方向，超过 5 年的预测没有意义。
\end{importantbox}

\subsection{本章小结}

选择 PhD 方向的关键不在于预测最准确，而在于\textbf{避免在信息已完全公开化的领域竞争}。三个可操作的原则：（1）不做现在最火的；（2）用``隔壁王阿姨''测试判断信息是否已过期；（3）保持 pivot 的灵活性，不做超过 5 年的锁定。


\section{未来方向：世界模型、能源与 AI 平民化}

\subsection{世界模型（World Model）：方兴未艾}

在讨论``赌什么方向''时，嘉宾给出了自己的判断：

\begin{quote}
以我有限的认知，计算机行业的话，我可能比较倾向于搞世界模型。还算是个方兴未艾吧，现在还没有太多公司在搞，效果也一般。Sora 2 也出了嘛。我可能觉得世界模型还可以再搞一波，搞个两三年，搞个 3 年左右。
\end{quote}

世界模型（World Model / Spatial Intelligence）在嘉宾看来处于典型的``已有人在做但远未成熟''的阶段——少数几家公司在探索，效果还远不能令人满意，但方向本身的空间是明确的。

\subsection{能源：非计算机的最佳赌注}

对于非计算机方向，能源是嘉宾的第一选择：

\begin{quote}
非计算机行业的话，那我只能选能源了。因为我开特斯拉太开心了，每天就感觉特斯拉这些电池、各种东西确实好。
\end{quote}

这个判断背后有更深层的逻辑：算力的爆炸式增长正在遭遇能源瓶颈。全球 60 多亿人，算力池要爆掉了，老黄（NVIDIA 黄仁勋）在卖卡，但电不够用。\textbf{数据中心的电力供应正在成为 AI 发展的硬约束}。

\subsection{AI 平民化：养猪场的启示}

讨论中最令人印象深刻的案例来自一位嘉宾的妻子所在的公司的经历：

\begin{quote}
有一个养猪场找到我老婆的公司，说``我想要检测这个猪的情绪''。这个东西在过去完全没有人问过。但这个老板用了 ChatGPT 之后，发现可以做这些事情，他们就有了很多不一样的想法。他想要检测猪的情绪和状态，因为这样可以去有效避免大规模猪瘟。
\end{quote}

\begin{importantbox}{AI 渗透正在引爆意想不到的领域}
第一次 AI 浪潮（CV 时代）只触达了少量科技公司。而 ChatGPT 的出现改变了这一点——它让那些从未想过用 AI 的传统行业从业者\textbf{自己产生了 AI 的想法}。养猪场老板主动找上门来要求 AI 服务，这是过去 CV 从业者做梦都想不到的需求场景。
\end{importantbox}

\subsection{传统行业的巨大信息差}

讨论延伸到传统行业与计算机行业的信息鸿沟：

\begin{itemize}
    \item 每个行业都希望用 AI，但每个行业都欠缺计算机基本原理的使用能力
    \item 非计算机、非互联网领域的信息化程度``极其落后''
    \item 这就意味着：程序员不需要在互联网这个狭窄领域卷，而是可以用自己的基础知识进入各种不同行业——仓储、物流、农业、法律
\end{itemize}

另一个案例：嘉宾的邻居是律师，律师行业现在也被要求学习使用 ChatGPT 来替代传统的人力整理档案和资料的工作，但他们学起来很吃力。

\subsection{多智能体系统：未来的架构师角色}

嘉宾提出了一个富有洞察力的``暴论''：

\begin{quote}
在未来一个 Multi-Agent 系统环境下，你最终处理的问题是一个类似于供应链的问题。你有一个 pipeline，这个 agent 产生了 bottleneck，agent 之间有反馈，你可以做各种在环节中的优化——比如应该提升这块的吞吐，优化那边的算法。那可能就会出现很多当年供应链专家，现在变成了所谓的``架构师''——产品架构师，去分析产品中的 bottleneck。
\end{quote}

\begin{knowledgebox}{多智能体系统架构师}
如果将 AI Agent 视为供应链中的节点，那么优化多 Agent 系统的运行效率——包括负载均衡、瓶颈检测、吞吐优化、反馈回路调优——将成为一个全新的专业领域。这个角色需要同时理解 AI 能力边界和系统设计原则。
\end{knowledgebox}

\subsection{本章小结}

未来 3 年的方向判断：世界模型（计算机）和能源（跨界）是两个有共识的赌注。但更宏观的机会在于 AI 向传统行业的渗透——这一波 AI 浪潮的特点是\textbf{需求端主动觉醒}，传统行业从业者首次主动寻求 AI 解决方案，而非被科技公司推销。多智能体系统的出现还将催生全新的架构师角色，这可能是 Infra 能力者的下一片蓝海。

\vspace{1em}
\begin{figure}[H]
\centering
\includegraphics[width=0.78\textwidth,keepaspectratio]{figures/talk_future.jpg}
\caption{讨论未来方向时的画面 \protect\footnotemark}
\end{figure}
\footnotetext{视频画面时间区间：00:32:00--00:36:00。}
